\documentclass[12pt, oneside]{book} 
%\usepackage[utf8]{vietnam}
%\usepackage[utf8]{english}
%\usepackage{dsfont}
\usepackage{amsmath, amssymb, color, graphicx, geometry, bbm}
\usepackage{float}
\usepackage{amsthm}
\usepackage{amsthm}
\usepackage{hyperref}
\usepackage{amsthm}
\usepackage{svg}
\usepackage{stmaryrd}
\usepackage{subfigure}
\usepackage{xcolor}
\usepackage{mathstyle}
\geometry{tmargin=.75in, bmargin=.75in, lmargin=.75in, rmargin = .75in}  

\newcommand{\C}{\mathbb{C}}
\newcommand{\Q}{\mathbb{Q}}
\newcommand{\Cdot}{\boldsymbol{\cdot}}

\newcommand{\nsubsection}[2]{
    \setcounter{subsubsection}{#1}
    \addtocounter{subsubsection}{-1}
    \subsubsection{#2}
}
\hypersetup{
    colorlinks = true
    linkcolor = true
}
\newcommand{\numberset}[1]{\llbracket #1\rrbracket}
%--------- \subsubsubsection ----------
\usepackage{titlesec}

\setcounter{secnumdepth}{4}

\titleformat{\paragraph}
{\normalfont\normalsize\bfseries}{\theparagraph}{1em}{}
\titlespacing*{\paragraph}
{0pt}{3.25ex plus 1ex minus .2ex}{1.5ex plus .2ex}
%--------- \subsubsubsection -----------

\newtheorem{thm}{Theorem}[chapter]
\newtheorem{defn}{Definition}[chapter]
\newtheorem{conv}{Convention}[chapter]
\newtheorem{rem}{Remark}[chapter]
\newtheorem{lem}{Lemma}[chapter]
\newtheorem{cor}{Corollary}[chapter]
\newtheorem{exercise}{Excercise}[section]
\numberwithin{equation}{chapter}

\title{An introduction to measure theory - T.Tao}
\author{Statistical Learning Theory Group}
\date{Academic Year 2026-2027}

\begin{document}

\maketitle
\tableofcontents

\vspace{.25in}

%\section{Introduction}

\chapter*{Preface}
\input{chapter0/0.0.1.tex}
\newpage
\begin{exercise}
    Let A, B be sets (maybe infinite), and $(x_{n,m})_{n\in A,m\in B},\quad x_{n,m}\in [0, +\infty]$ be a doubly infinite sequence of extended non-negative reals. Show that $$\sum_{(n,m) \in A \times B} x_{n,m} = \sum_{n\in A}\sum_{m\in B} x_{n,m} = \sum_{m\in B}\sum_{n\in A} x_{n,m}$$
\end{exercise}
\textbf{Intuition:}
\begin{itemize}
    \item For tenary equation, we will prove 2 binary equalities then bridge them. We can use same strategy for Tonelli's theorem for reference. 
    \item Cartesian over finite sets can be written down to collections of collections.
    \item Tonelli's theorem: nonnegativity + directed finite approximations
\end{itemize}
\begin{proof}
    Let $F \subseteq A \times B$ be a finite subset of Cartesian product of A and B. We define $F_A = \{n \in A: \exists m \in B, (n, m) \in F\}$ and $F_B = \{m \in B: \exists n \in A, (n, m) \in F\}$. It suffices to show that $F_A$ and $F_B$ are finite sets, and $F_A \subseteq A, F_B \subseteq B$. F is then can be written as subset of Cartesian of finite sets (big brain move) $F \subseteq F_A \times F_B$. As a result, we have:
    \[
        \sum_{(n,m) \in F} x_{n,m} \leq \sum_{(n, m) \in F_A \times F_B} x_{n,m} = \sum_{n\in F_A}\sum_{m\in F_B} x_{n,m} \leq \sum_{n\in A}\sum_{m\in B} x_{n,m} \quad (a)
    \]
    We need to prove that:     
    \[
    \sup_{F}  \sum_{(n,m) \in F} x_{n,m} = \sup_{F_A, F_B} \sum_{n\in F_A}\sum_{m\in F_B} x_{n,m}\qquad \text{(1)} 
    \]
    \[
    \sup_{F_A, F_B} \sum_{n\in F_A}\sum_{m\in F_B} x_{n,m} = \sum_{n \in A}\sum_{m \in B} x_{n,m}\qquad \text{(2)} 
    \]
    \textbf{Lemma 1:} If N is finite and $a_{n, M} = \sum_{m \in M} x_{n,m} \in [0, +\infty]\quad \forall n\in N$, then $\sum_{n\in N} sup_{M} a_{n,M} = sup_{M} \sum_{n\in N} a_{n,M}$ \\
    \textbf{Proof of Lemma 1:} For every M, $a_{n, M} \leq \sup_M a_{n,M}$ (definition of supremum), therefore $\sup_M \sum_{n\in N} a_{n, M} \leq \sum_{n\in N} sup_M a_{n,M}$. On the contrary, we need to prove $\sum_{n\in N}  \sup_M a_{n,M} \leq\sup_M \sum_{n\in N} a_{n, M}$ to establish equality. This proof has to solve two cases: (1) if $\sum_{n\in N} \sup_M a_{n, M} = \infty$ then there exists at least one element $a_{n_i, M_i} = \infty$ since $N$ is finite, therefore $sup_{M} \sum_{n\in N} a_{n,M} = \infty$ which proves the point, (2) if $a_{n,M}$ is real then M is finite by constraint, by definition of supremum, we can always choose a finite subset $M_n \subset M$ such that $\sup_{M} \sum_{m\in M} a_{n, M} - \frac \epsilon ||N|| \leq a_{n, M_n} \leq \sup_{M} \sum_{m\in M} a_{n, M}$, where gathered sum induces: 
    \[
    \sum_{n\in N} a_{n,M} \geq \sum_{n\in N} \sup_M a_{n,M} - \epsilon
    \]
    This is true for every $\epsilon > 0$, therefore $\sup_{M} \sum_{n\in N} a_{n,M} \geq \sum_{n\in N} \sup_M a_{n,M}$.

    For (2), this needs to be explicitly proven through induction. For each fixed n, we can see that: 
    \[
    \sum_{m\in B} x_{n,m} = \sup_{M' \subseteq B, M' \text{ finite}} \sum_{m \in M'} x_{n,m}
    \]
    Pack this sum as element and with A, we results in final supremum of finite sum:
    \[
        \sum_{n\in A}\sum_{m\in B} x_{n,m} = \sup_{\substack{F_A \subseteq A, F_A \text{ finite} \\ F_B \subseteq B, F_B \text{ finite}}} \sum_{n\in F_A}\sum_{m\in F_B} x_{n,m} 
    \]


    For (1), denoting left and right sides of (1) as $S_1$ and $S_2$ respectively, (a) has proven that $\sum_{(n,m) \in F} x_{n,m} \leq \sum_{(n, m) \in F_A \times F_B} x_{n,m} \leq S_2$ for every finite set of F, therefore $S_1 \leq S_2$ (1a). On the other hand,  $\sum_{(n, m) \in F_A \times F_B} x_{n,m} = \sum_{n\in F_A}\sum_{m\in F_B} x_{n,m}$. For $F_A \subseteq A$ and $F_B \subseteq B$, we would have $F_A \times F_B \subseteq A \times B$ whose supremum is taken by $S_1$, which results in $\sum_{n\in F_A}\sum_{m\in F_B} x_{n,m} \leq S_1$ for every finite $F_A$ and $F_B$. Therefore, $S_2 \leq S_1$ (1b). As a result from (1a) and (1b), equation (1) is proven.

    In conclusion, combining (1) and (2), we will have: 
    \[
    \sum_{(n,m) \in A \times B} x_{n,m} = \sup_{F}  \sum_{(n,m) \in F} x_{n,m} = \sup_{F_A, F_B} \sum_{n\in F_A}\sum_{m\in F_B} x_{n,m} = \sum_{n \in A}\sum_{m \in B} x_{n,m}
    \]
    The same proof goes for equation: $ \sum_{(n,m) \in A \times B} x_{n,m} = \sum_{m \in B}\sum_{n \in A} x_{n,m}$
\end{proof}

\chapter{Measure Theory}
\section{Prologue: The problem of measure}
\input{chapter1/1.1.1.tex}
\textbf{Lemma 1.1.2} Let $E \subset R^d$ be an elementary set.
(ii) If $E = \cup_{i=1}^n B_i$, and $E = \cup_{j=1}^m C_j$
Then $\sum_{i=1}^n |B_i| = \sum_{j=1}^m |B_j|$

\textbf{Proof}

First, we need prove ${B_i \cap C_j}_{i,j}$ are boxes. Considering
\[
\begin{split}
B_i = I_{i,1} \times \dots \times I_{i,d}\\
C_j = J_{j,1} \times \dots \times J_{j,d}\\
I_{i,r} \cap J_{j,r} = I' || \emptyset
\end{split}
\]
which implies $B_i \cap C_j = box || \emptyset $

Suppose overlap 
\[
(B_i \cap C_j) \cap (B_i' \cap C_j') \neq \emptyset
\]
\[
\implies x \in B_i \cap B_i' \text{ and } x \in C_j \cap C_j'
\]
But $B_i, B_i'$ disjoint which cause a contradiction.
Hence, $B_i \cap C_j$ mutually disjoint.

$\forall x \in E$, always $\exists i$ that $x \in B_i$ and j that $x \in C_j$
This means $x \in B_i \cap C_j \implies E \subseteq \cup_{i,j}(B_i \cap C_j)$ 
We also have
\[
B_i \cap C_j \subseteq E \implies \cup_{i,j}(B_i \cap C_j) \subseteq E
\]
Combine these equations , we have
\[
E = \sqcup_{i,j} (B_i \cap C_j)
\]
Apply measure function and sum up with double indexing result the desired equality.

\input{chapter1/1.1.3.tex}
\input{chapter1/1.1.4.tex}

\input{chapter1/1.1.7.tex}
\input{chapter1/1.1.8.tex}
\input{chapter1/1.1.9.tex}


\input{chapter1/1.1.11.tex}


\noindent\textbf{Exercise 1.1.14 (Dyadic approximation and Jordan measurability).}\par
Let $E\subset\mathbb{R}^d$ be bounded. For $n\geq 0$, let
$\mathcal D_n$ be the dyadic cubes of side length $h_n=2^{-n}$,
and set
\[
N_n^-:=\#\{Q\in\mathcal D_n:Q\subseteq E\},\qquad
N_n^+:=\#\{Q\in\mathcal D_n:Q\cap E\neq\varnothing\}.
\]
Show that $E$ is Jordan measurable if and only if
\[
2^{-dn}(N_n^+-N_n^-)\longrightarrow0.
\]
In this case, show that both $2^{-dn}N_n^-$ and $2^{-dn}N_n^+$
converge to $m(E)$.

\begin{proof}
For any bounded $T\subset\mathbb R^d$, write
\[
I_n(T):=\bigcup_{\substack{Q\in\mathcal D_n\\Q\subseteq T}}Q,
\qquad
O_n(T):=\bigcup_{\substack{Q\in\mathcal D_n\\Q\cap T\neq\varnothing}}Q.
\]
Only finitely many cubes at scale $n$ intersect $T$; hence
$I_n(T)$ and $O_n(T)$ are elementary. The cubes at each scale
partition $\mathbb R^d$, so
\[
I_n(T)\subseteq T\subseteq O_n(T),\qquad
m(I_n(E))=2^{-dn}N_n^-,\quad m(O_n(E))=2^{-dn}N_n^+.
\tag{1}
\]

\emph{Claim.} If $S$ is elementary, then
\[
m(I_n(S))\longrightarrow m(S),\qquad
m(O_n(S))\longrightarrow m(S).
\tag{2}
\]
Indeed, for an interval $J\subset\mathbb R$ of length $\ell$,
only the dyadic intervals meeting its two endpoints can contribute
to the difference between inner and outer approximations. Thus,
with $h=h_n$,
\[
(\ell-2h)_+\leq m_1(I_n(J))\leq\ell
\leq m_1(O_n(J))\leq\ell+2h.
\]
Consequently, both one-dimensional approximation measures tend
to $\ell$. For a box $R=J_1\times\cdots\times J_d$, the dyadic
partition is a product partition, so
\[
I_n(R)=\prod_{j=1}^d I_n(J_j),\qquad
O_n(R)=\prod_{j=1}^d O_n(J_j).
\]
The elementary product formula now gives
$m(I_n(R))\to |R|$ and $m(O_n(R))\to |R|$.
For an elementary set $S=\bigsqcup_{k=1}^M R_k$, note that
\[
\sum_{k=1}^M m(I_n(R_k))\leq m(I_n(S))\leq m(S)
\leq m(O_n(S))\leq\sum_{k=1}^M m(O_n(R_k)).
\]
Both outermost expressions converge to
$\sum_{k=1}^M|R_k|=m(S)$, proving (2) by squeezing.

Suppose first that $E$ is Jordan measurable. Given
$\varepsilon>0$, Exercise~1.1.5 supplies elementary sets
$A\subseteq E\subseteq B$ such that
$m(B\setminus A)<\varepsilon/3$.
Since dyadic inner and outer approximations preserve inclusions,
\[
I_n(A)\subseteq I_n(E)\subseteq O_n(E)\subseteq O_n(B).
\]
By (2), for all sufficiently large $n$,
\[
0\leq m(A)-m(I_n(A))<\varepsilon/3,
\]
and
\[
0\leq m(O_n(B))-m(B)<\varepsilon/3.
\]
Therefore
\[
\begin{aligned}
0\leq m(O_n(E))-m(I_n(E))
&\leq m(O_n(B))-m(I_n(A))\\
&=[m(O_n(B))-m(B)]+[m(B)-m(A)]\\
&\qquad+[m(A)-m(I_n(A))] < \varepsilon.
\end{aligned}
\]
By (1), $2^{-dn}(N_n^+-N_n^-)\to0$.

Conversely, suppose $2^{-dn}(N_n^+-N_n^-)\to0$.
For every $\varepsilon>0$ we may choose $n$ with
\[
m(O_n(E)\setminus I_n(E))
=m(O_n(E))-m(I_n(E))<\varepsilon.
\]
Since $I_n(E)\subseteq E\subseteq O_n(E)$ and both approximations
are elementary, Exercise~1.1.5 implies that $E$ is Jordan measurable.

Finally, for a Jordan measurable $E$, monotonicity gives
\[
m(I_n(E))\leq m(E)\leq m(O_n(E)).
\]
The difference of the outer bounds tends to zero, so both
converge to $m(E)$. Apply (1) to obtain
\[
\displaystyle
m(E)=\lim_{n\to\infty}2^{-dn}N_n^-
=\lim_{n\to\infty}2^{-dn}N_n^+.
\]
\end{proof}

\end{document}
